\begin{proof}
\textbf{Young’s Inequality.}  
Let \(a,b\ge0\) and let \(p,q>1\) satisfy \(\dfrac1p+\dfrac1q=1\).  
We prove  

\[
ab\le\frac{a^{p}}{p}+\frac{b^{q}}{q},
\]

with equality exactly when \(a^{p}=b^{q}\).

\begin{enumerate}
\item \textbf{Conjugate exponents.}  
From \(\dfrac1p+\dfrac1q=1\) we obtain the familiar identities  

\[
q=\frac{p}{p-1},\qquad p=\frac{q}{q-1}. \tag{C}
\]

\item \textbf{Degenerate cases.}  
If \(a=0\) or \(b=0\) then both sides of the desired inequality are zero, so the claim holds (with equality).  
Henceforth assume  

\[
a>0,\qquad b>0. \tag{D}
\]

\item \textbf{Homogenisation.}  
Dividing the inequality by the positive product \(ab\) yields  

\[
1\le\frac{a^{p-1}}{p\,b}+\frac{b^{q-1}}{q\,a}. \tag{1}
\]

Introduce  

\[
u:=\frac{a^{p-1}}{b},\qquad v:=\frac{b^{q-1}}{a}, \tag{2}
\]

so that (1) becomes  

\[
\frac{u}{p}+\frac{v}{q}\ge1. \tag{3}
\]

\item \textbf{Weighted AM–GM.}  
The weighted arithmetic–geometric mean inequality states that for non‑negative \(x,y\) and weights \(\lambda,\mu\ge0\) with \(\lambda+\mu=1\),

\[
\lambda x+\mu y\ge x^{\lambda}y^{\mu},
\]

with equality iff \(x=y\) (or a weight equals \(1\)).  
Choose  

\[
\lambda=\frac1p,\qquad \mu=\frac1q
\quad(\text{so }\lambda+\mu=1\text{ by }(C)),
\]

and set \(x=u,\;y=v\).  Then

\[
\frac{u}{p}+\frac{v}{q}\ge u^{1/p}v^{1/q}. \tag{4}
\]

\item \textbf{Simplification of the right–hand side.}  
Using the definitions (2),

\[
\begin{aligned}
u^{1/p}v^{1/q}
&=\left(\frac{a^{p-1}}{b}\right)^{\!1/p}
  \left(\frac{b^{q-1}}{a}\right)^{\!1/q} \\[2mm]
&=a^{\frac{p-1}{p}}\,b^{-\frac1p}\,
  b^{\frac{q-1}{q}}\,a^{-\frac1q}.
\end{aligned}
\]

Because \(\frac1p+\frac1q=1\) (equivalently (C)), we have  

\[
a^{\frac{p-1}{p}-\frac1q}=a^{1-\frac1p-\frac1q}=a^{0}=1,
\qquad
b^{\frac{q-1}{q}-\frac1p}=b^{0}=1,
\]

hence  

\[
u^{1/p}v^{1/q}=1. \tag{5}
\]

\item \textbf{Conclusion of the core inequality.}  
Combining (4) and (5) gives  

\[
\frac{u}{p}+\frac{v}{q}\ge1,
\]

which is exactly (3). Restoring \(u,v\) from (2) we obtain (1).

\item \textbf{Returning to the original variables.}  
Multiplying (1) by the positive quantity \(ab\) (permitted by (D)) yields  

\[
ab\le\frac{a^{p}}{p}+\frac{b^{q}}{q},
\]

the desired Young’s inequality.

\item \textbf{Equality condition.}  
Equality in the weighted AM–GM step (4) occurs iff \(u=v\).  From (2),

\[
u=v\;\Longleftrightarrow\;
\frac{a^{p-1}}{b}=\frac{b^{q-1}}{a}
\;\Longleftrightarrow\;
a^{p}=b^{q}.
\]

Together with the degenerate case \(a=0\) or \(b=0\) (where both sides vanish), the complete equality condition is:

\begin{itemize}
\item either \(a=0\) or \(b=0\);
\item or \(a,b>0\) and \(a^{p}=b^{q}\).
\end{itemize}
Thus the statement is proved. ∎
\end{enumerate}
\end{proof}